\chapter*{Abstract}
\addcontentsline{toc}{chapter}{Abstract}

This paper presents an advanced deep learning methodology for automated corpus callosum segmentation in brain MRI images. Leveraging an **Enhanced UNet** architecture, our approach integrates **residual connections**, **attention mechanisms**, and **Atrous Spatial Pyramid Pooling (ASPP)** to enhance feature extraction and context integration. Furthermore, **deep supervision** is incorporated during training to improve convergence and feature learning across multiple scales. We evaluate our model on a comprehensive dataset of 1,450 T1-weighted sagittal MRI scans, achieving a **Dice Similarity Coefficient of 0.923 ± 0.024** and an **Intersection over Union (IoU) of 0.858 ± 0.041**. Ablation studies confirm the significant contribution of each architectural enhancement. Our model consistently outperforms state-of-the-art segmentation methods, demonstrates high computational efficiency, and shows excellent generalization across diverse patient populations. This automated solution offers substantial improvements in consistency and speed compared to manual segmentation, holding significant clinical implications for neuroimaging research and diagnosis.

\textbf{Keywords:} Corpus Callosum, Segmentation, Deep Learning, UNet, Attention, ASPP, Deep Supervision, MRI, Brain Imaging, Dice Score